\begin{exsol}{5}{\MR{} Food for Thought 4.1}{119}{Foundations}{A1}
\asked Explain whether efficiency can justify the switch of market economies to central planning
of scarce resources in the two world wars, and what other (non-efficiency) reasons support it.

\method The three criteria for comparing coordination systems: efficiency with perfect
information, communication, brittleness (\MR{} Ch.~4, pp.~93--94). Prices versus quantities and
Equation~4.1 (pp.~94--98). Design and innovation attributes (pp.~91--93) and planning with design
attributes (pp.~103--106). The failure of prices under scale economies (summary, p.~117). \MR{}
note that planners in the most market-oriented economies directed steel, rubber, sugar and meat
in the Second World War (p.~89).

\solution
\subparagraph{Efficiency reasons.}
\begin{enumerate}[label=(\roman*)]
  \item \emph{A sudden, huge change with innovation attributes.} War replaces the peacetime
    pattern of output with a new one that nobody in the market has experience of: ships, aircraft,
    munitions in enormous quantities. The information on what is needed sits with the military
    and the government, not with dispersed buyers and sellers. A price system would have to grope
    its way to an entirely new equilibrium (p.~103), and this takes time a war does not allow.
    Problems whose optimum depends on information that operating units lack are badly solved by
    decentralized prices (p.~92).
  \item \emph{Design attributes on a national scale.} War production is a giant synchronization
    and fit problem. Aircraft need engines, aluminium, fuel, trained crews and shipping at the
    same time and place. An operation needs everything in the same week. A priori the pattern of
    the solution is known (the parts must match), and a misfit is far more costly than a
    slightly wrong overall scale (p.~91, p.~105). Announcing quantities and dates is then both
    cheaper in communication and less brittle than prices (pp.~104--106).
  \item \emph{Brittleness.} By Equation~4.1 (p.~97), the loss under a price control relative to a
    quantity control equals the slope of marginal benefit over the slope of marginal cost. For
    critical war goods the marginal benefit is very steep: falling short of the required quantity
    of ammunition or fuel is disastrous, and a surplus is worth little. Marginal cost, by contrast,
    is fairly flat once mass production is set up. A steep marginal benefit relative to marginal
    cost is exactly the case in which a quantity directive loses less than a price when the
    planner's estimates are wrong (pp.~96--98).
  \item \emph{No market to price against.} The government is the single buyer of munitions, and
    the new mass-production lines have large scale economies, so market-clearing prices may not
    exist or may not guide output efficiently (p.~117).
\end{enumerate}

\subparagraph{Other reasons.}
\begin{itemize}
  \item \emph{Distribution and morale.} With civilian goods cut back, market prices would ration
    sugar or meat by ability to pay. Coupon rationing spreads the sacrifice evenly, which matters
    for consent to the war. This is a distributional argument, not an efficiency one.
  \item \emph{Inflation and windfall gains.} A sharp rise in prices transfers large rents to the
    owners of scarce inputs and raises the government's own procurement bill. Price controls with
    rationing suppress both, though at the cost of shortages and black markets.
  \item \emph{Urgency and authority.} War concentrates decision power and gives the government
    a single, widely shared objective, which reduces the political cost of command.
\end{itemize}

\subparagraph{The other side.} Central planning keeps its usual costs: it ignores local
information about costs, it invites the motivation problems of planned economies (the ratchet
effect, p.~14), and it breeds black markets. That is why even wartime economies planned only the
critical sectors, left many decisions to firms and consumers, and dismantled controls after the
war. The efficiency case is strongest for the military-industrial core (design and innovation
attributes, steep marginal benefit) and weakest for consumer goods, where the case for rationing
is mainly distributional.

\interp Which coordination system is best depends on the problem, not on ideology. The same
criteria that favour prices for routine resource allocation (Hurwicz, pp.~100--103) favour
quantities and central direction when a problem has design and innovation attributes, urgency and
a steep cost of shortfall.

\verify Consistency check with Table~4.1 (p.~95). There, with the planner wrong about costs, a
quantity directive loses $1$ unit of net benefit and a price directive loses $15$. The table is
the case of a marginal benefit that is steep relative to marginal cost --- the wartime case. In
peacetime, with flat marginal benefit (many substitutes) and steep marginal cost, Eq.~4.1 reverses
the ranking, consistent with returning to prices after the war.

\answer{Yes. War creates an abrupt change with innovation attributes, a national problem with
design attributes (synchronization and fit of production), a steep marginal benefit of critical
goods that makes quantity directives less brittle than prices (Eq.~4.1), and monopsony and scale
economies that undermine price guidance. Non-efficiency reasons --- equitable rationing for
morale, curbing inflation and windfall profits, urgency --- reinforce the case, which is strongest
for military production and mainly distributional for consumer goods.}
\end{exsol}
